% SOURCE_AUTHORITY: sga5_fr_workpass.tex, lines 8965--8978 (LF-normalized unit).
% SOURCE_SHA256: 791F4EFFC5E02832D5D77ED03518C8156D6F07E4C8238B03545DB93D883FBB28
\subsection*{3. Classe de cohomologia \(\ell\)-ádica associada a um ciclo}

\subsubsection*{3.1}
Se \(X\) é um esquema próprio, ou liso e de tipo finito, sobre um corpo algebricamente fechado \(k\), sabe-se que, para todo feixe construtível \(M\), abeliano e de torção prima com a característica de \(k\), os grupos de cohomologia
\[
H^i(X,M)
\]
São grupos abelianos de comprimento finito. Isso permite (2.2.1, Observação) definir, para todo \(\mathbb Z_\ell\)-feixe \(F\) (\(\ell\) primo com a característica de \(k\)), um sistema projetivo \((AR,\ell)\)-ádico noetheriano de grupos abelianos
\[
H^i(X,F)
\]
com as propriedades habituais. Em particular, o método utilizado no número 2 permite definir o produto cup, o homomorfismo de imagem inversa para um morfismo de esquemas sobre \(k\) e o homomorfismo de imagem direta para um morfismo próprio de esquemas suaves sobre \(k\) (cf. Exp. IV para as definições em relação aos feixes de torção).

Por outro lado, sabe-se (1.2.4) que a categoria dos \(\mathbb Z_\ell\)-feixes sobre um corpo algebricamente fechado é equivalente à categoria dos \(\mathbb Z_\ell\)-módulos de tipo finito, e que a equivalência é dada pelo funtor limite projetivo. Na prática, será conveniente identificar o \(\mathbb Z_\ell\)-módulo de tipo finito \(H^i(X,F)\) com o \(\mathbb Z_\ell\)-feixe correspondente.
